\lecture{5}
\title{Fields, Vector Spaces, Dimension}

\begin{document}

So far we've directed our attention to properties and structures relating to groups.  In this lecture, we'll redirect our attention toward the applications of all these properties and structures.

\subsection{Field Definitions}

\textbf{Definition (Field).} A \textbf{field} $F$ is a set $F$ together with \emph{two} operations $+$ and $\cdot$ satisfying the following properties:
\begin{enumerate}
    \item The set $(F, +)$ is an abelian group with identity $0 \in F$.
    \item The set $(F \setminus \{0\}, \cdot)$ is an abelian group with identity $1 \in F$.
    \item For any $x, y, z \in F$, the identity $x \cdot (y + z) = x \cdot y + x \cdot z$ holds. \emph{(Distributive Law)}
\end{enumerate}

\begin{remark}
    All fields $F$ contain an image of $\mathbb{Z}$; that is, it makes sense to refer to ``$5$'' or ``$-67$'' as elements of $F$.  This is because $1$ is an element of $F$, so when we refer to ``$5$'' in $F$, we're just referring to ``$1 + 1 + 1 + 1 + 1$''.
\end{remark}

\textbf{Example.} All of $(\mathbb{R}, +, \cdot)$, $(\mathbb{C}, +, \cdot)$, and $(\mathbb{Q}, +, \cdot)$ are fields.  Meanwhile, $(\mathbb{Z}, +, \cdot)$ is not a field.

Most of the algebraic properties we expect from $(\mathbb{R}, +, \cdot)$ hold in general fields.

\begin{theorem}{Additive Identity is Annihilating}
    For all $x \in F$, the identity $x \cdot 0 = 0$ holds.
\end{theorem}

\begin{proof}
    Since $0$ is the additive identity, $y + 0 = y$.  Multiplying by $x$ on both sides yields $x \cdot y + x \cdot 0 = x \cdot y$.  Subtracting $x \cdot y$ from both sides yields $x \cdot 0 = 0$.
\end{proof}

\begin{theorem}{Zero Product Property}
    For any $x, y \in F$, if $xy = 0$, then either $x = 0$ or $y = 0$.
\end{theorem}

\begin{proof}
    If $x \neq 0$, then $x$ has a multiplicative inverse.  Multiplying by $x^{-1}$ on both sides of $xy = 0$ yields $y = x^{-1} \cdot 0 = 0$.  Thus, if $x \neq 0$, then $y = 0$, and we're done. ~ $\blacksquare$
\end{proof}

Of particular note are \emph{finite} fields.  We can name some in the theorem below.

\begin{theorem}{Finite Fields}
    The set $(\mathbb{Z}/n \mathbb{Z}, +, \cdot)$ is a field if and only if $n$ is prime.
\end{theorem}

\begin{proof}
    The only property of $(\mathbb{Z}/n\mathbb{Z}, +, \cdot)$ we need to check is that multiplicative inverses exist.

    If $n = ab$ is composite, then $a$ does not have a multiplicative inverse, so $(\mathbb{Z}/n\mathbb{Z}, +, \cdot)$ is not a field.

    If $n$ is prime, then we claim all nonzero $a \in \mathbb{Z}/n\mathbb{Z}$ have a multiplicative inverse.  Here are two proofs:

    \underline{Proof \#1:} Consider the additive subgroup $\langle a \rangle$ of $\mathbb{Z}/n\mathbb{Z}$.  By Lagrange's Theorem, the order of this subgroup divides $n$.  But $n$ is prime, so the order of this subgroup must either be $1$ or $n$.
    
    The former case would mean $\langle a \rangle = \{0\}$, i.e. $a = 0$, which we've excluded.  So we're in the latter case, where $\langle a \rangle = \mathbb{Z}/n\mathbb{Z}$.  Then $\langle a \rangle$ contains $1$, meaning there is some integer $b \in \mathbb{Z}$ such that $ab = 1$, as desired. ~ $\square$

    \underline{Proof \#2:} Equivalently, it suffices to show that for any $a \in \mathbb{Z}$ such that $\gcd(a, n) = 1$, there must exist $b, c \in \mathbb{Z}$ such that $ab + nc = 1$.  The key idea is to consider the group $a\mathbb{Z} + n\mathbb{Z}$.  It must be a subgroup of $\mathbb{Z}$, so it looks like $k\mathbb{Z}$ for some $k \in \mathbb{Z}$.  This $k$ must satisfy the property that $k \mid a$ and $k \mid n$.  But $\gcd(a, n) = 1$, so it must be that $k = 1$.
    
    In other words, $a\mathbb{Z} + n\mathbb{Z} = \mathbb{Z}$.  So there exist some $b, c \in \mathbb{Z}$ such that $ab + nc = 1$, exactly as desired. ~ $\square$
\end{proof}

Thus, the set $(\mathbb{Z}/n\mathbb{Z}, +, \cdot)$ is a field for primes $n$.  But be careful!  This does not hold when $n$ is a higher power of a prime, since (e.g.) $3 \in (\mathbb{Z}/9\mathbb{Z}, +, \cdot)$ has no multiplicative inverse.

\begin{remark}
    There still do exist finite fields with prime-power order, though.  They just look different.
\end{remark}

\subsection{Vector Space Definitions}

\textbf{Definition (Vector Space).} Let $F$ be a field.  Call the elements of $F$ ``scalars''.  Then $V$ is a \textbf{vector space over $\mathbf{F}$} if:
\begin{itemize}
    \item It has two binary operations: abelian-group vector addition $+: V \times V \to V$ and scalar multiplication $\cdot: F \times V \to V$.
    \item The multiplicative identity $1 \in F$ satisfies $1v = v$ for all $v \in V$.
    \item For any $\lambda, \mu \in F$ and any $v \in V$, we have $(\lambda \mu) \cdot v = \lambda(\mu \cdot v)$. \emph{(Associative Law)}
    \item For any $\lambda, \mu \in F$ and any $u, v \in V$, we have $\lambda(u + v) = \lambda u + \lambda v$ and $(\lambda + \mu)v = \lambda v + \mu v$. \emph{(Distributive Law)}
\end{itemize}
Note, importantly, that there is no sense of vector multiplication.

\textbf{Example.} Naturally, $V = \mathbb{R}^n$ is a vector space over $\mathbb{R}$.  Some less obvious examples include:
\begin{itemize}
    \item $\mathbb{C}$ is a two-dimensional vector space over $\mathbb{R}$.
    \item $\mathbb{R}$ is an infinite-dimensional vector space over $\mathbb{Q}$.
    \item $\{f: \mathbb{R} \to \mathbb{R}\}$ is an infinite-dimensional vector space over $\mathbb{R}$.
    \item $\{y: \mathbb{R} \to \mathbb{C}: y'' = -y\}$ is a two-dimensional vector space over $\mathbb{C}$, or a four-dimensional vector space over $\mathbb{R}$.
\end{itemize}

Recall that we can define matrices $M \in \mathbb{R}^{m \times n}$ as linear maps $M: \mathbb{R}^n \to \mathbb{R}^m$ between vector spaces $\mathbb{R}^m$ and $\mathbb{R}^n$ over $\mathbb{R}$.  Let's generalize this so that it applies to general fields $F$ beyond just $\mathbb{R}$.

\textbf{Definition (Matrices).} Consider vector spaces $F^m$ and $F^n$ over a field $F$.  Then a matrix $M \in F^{m \times n}$ is both (i) a rectangular array consisting of elements of $F$, and (ii) a linear map from vectors in $F^n$ to vectors in $F^m$.
\[ M = \begin{bmatrix} M_{11} & M_{12} & \cdots & M_{1n} \\ M_{21} & M_{22} & \cdots & M_{2n} \\ \vdots & \vdots & \ddots & \vdots \\ M_{m1} & M_{m2} & \cdots & M_{mn} \end{bmatrix} ~ \implies ~ M\begin{bmatrix} v_1 \\ v_2 \\ \vdots \\ v_n \end{bmatrix} = \begin{bmatrix} M_{11}v_1 + M_{12}v_2 + \cdots + M_{1n}v_n \\ M_{21} v_1 + M_{22} v_2 + \cdots + M_{2n} v_n \\ \vdots \\ M_{m1} v_1 + M_{m2}v_2 + \cdots + M_{mn} v_n \end{bmatrix} \]
In other words, we just extend the notational definitions from matrices in $\mathbb{R}^{m \times n}$ verbatim to matrices in $F^{m \times n}$.  Furthermore, all the row-reduction properties of matrices in $\mathbb{R}^{m \times n}$ work for matrices in $F^{m \times n}$, too.

\begin{remark}
    It does not make sense to describe linear maps between \emph{arbitrary} vector spaces $M: V \to W$ using a matrix.  Matrices only make sense if the vector spaces $V$ and $W$ are both of the form $F^m$ and $F^n$.
\end{remark}

\textbf{Definition (Span).} Let $V$ be a vector space over $F$.  Then for any set of vectors $\{v_1, \dots, v_n\} \subseteq V$, their \textbf{span} is:
\[ \mathrm{span}\left(\{v_1, \dots, v_n\}\right) = \left \{ \sum_{i = 1}^n \lambda_i v_i: \lambda_i \in F \right \}. \]
More generally, for any potentially infinite set $S \subseteq V$ of vectors, the \textbf{span} of $S$ is:
\[ \mathrm{span}\left(S\right) = \{\text{linear combinations of \emph{finitely} many elements of } S\}. \]
Infinite sums in general vector spaces are nonsensical because there is no inherent notion of convergence.

\begin{remark}
    It is equivalent to say that the span of $S \subseteq V$ is the smallest vector subspace of $V$ containing $S$.
\end{remark}

\textbf{Definition (Spanning).} We say that a set $S \subseteq V$ of vectors \textbf{spans} $V$ if $\mathrm{span}(S) = V$.

\textbf{Definition (Linear Independence).} A set $S \subseteq V$ of vectors is \textbf{linearly independent} if any finite linear combination of vectors in $S$ equals zero if and only if all coefficients are zero.  In other words,
\[ \text{For all distinct } v_1, \dots, v_n \in S, \text{ we have } \left ( \sum_{i = 1}^n \lambda_i v_i = 0 \right) \iff \left( \lambda_i = 0 \text{ for all } i \right)  \]

Here's an equivalent definition of linear independence.

\begin{theorem}{Linear Independence vs. Unique Linear Combinations}
    A set $S \subseteq V$ of vectors is linearly independent if and only if every element in $\mathrm{span}(S)$ has a \emph{unique} expression as a linear combination of $S$.
\end{theorem}

\begin{proof}
    Argue by contradiction: suppose some $v \in \mathrm{span}(S)$ has two distinct linear combinations of $S = \{v_1, \dots, v_n\}$.
    \[ \begin{cases} v = \lambda_1v_1 + \dots + \lambda_n v_n \\
    v = \mu_1 v_1 + \dots + \mu_n v_n \end{cases}. \]
    Subtracting these two linear combinations yields $0 = (\lambda_1 - \mu_1)v_1 + \dots + (\lambda_n - \mu_n)v_n$, which is a nontrivial linear combination of $S$ equal to zero.  So $S$ is not linearly independent, contradiction.
    
    Conversely, if $S$ is not linearly independent, then some nontrivial linear combination of $S$ equals zero, giving $0$ a second expression alongside the all-zero one. ~ $\blacksquare$
\end{proof}

\subsection{Basis and Dimension}

Now that we've defined span and linear independence, we can talk about bases.

\textbf{Definition (Basis).} A set $S \subseteq V$ of vectors is a \textbf{basis} of $V$ if it both spans $V$ and is linearly independent.

\textbf{Example.} One basis of $\mathbb{R}^3$ is $\left\{ \begin{bsmallmatrix} 1 \\ 0 \\ 0 \end{bsmallmatrix}, \begin{bsmallmatrix} 0 \\ 1 \\ 0 \end{bsmallmatrix}, \begin{bsmallmatrix} 0 \\ 0 \\ 1 \end{bsmallmatrix} \right \}$.  Another equally valid basis is $\left\{ \begin{bsmallmatrix} 1 \\ 0 \\ 1 \end{bsmallmatrix}, \begin{bsmallmatrix} 2 \\ 0 \\ 0 \end{bsmallmatrix}, \begin{bsmallmatrix} 1 \\ 3 \\ 5 \end{bsmallmatrix} \right \}$.

\textbf{Example.} Importantly, vector spaces need not have a ``natural'' or ``inherently nice'' basis.
\begin{itemize}
    \item One basis of $\{y: \mathbb{R} \to \mathbb{C}: y'' = -y\}$ is $\{ \cos x, \sin x\}$.  Another basis is $\{e^{ix}, e^{-ix}\}$.  There is no way to determine which of these two bases is better (or whether, say, $\{\cos x + e^{ix}, \sin x + e^{-ix}\}$ is a better basis, either!).
    \item There are many bases of the vector space $\mathbb{R}$ over $\mathbb{Q}$, but they're impossible to describe explicitly.
\end{itemize}

We want to use the size of a basis of $V$ to talk about the dimension of $V$.  Here's a first step.

\textbf{Definition (Finite Dimension).} A vector space $V$ is said to be \textbf{finite-dimensional} if and only if it has a finite spanning set (and therefore a finite basis).

\begin{theorem}{Existence of Basis}
    Every finite-dimensional vector space $V$ has a basis $S \subseteq V$.
\end{theorem}

\begin{proof}
    Since $V$ is finite-dimensional, there exists a finite set $S = \{v_1, \dots, v_n\} \subseteq V$ of vectors that spans $V$.  If $S$ is also linearly independent, then we're done: $S$ is our basis.  If $S$ is not linearly independent, that means:
    \[ \lambda_1 v_1 + \dots + \lambda_n v_n = 0 \text{ for some } \{\lambda_i\} \in F \text{ not all zero.} \]
    Pick some index $i$ such that the coefficient $\lambda_i$ is nonzero.  We claim that for this $i$, we can just remove $v_i$ from $S$ and still be left with a spanning set $S' := S \setminus \{v_i\}$.  The reason why is the following:
    \[ \lambda_1 v_1 + \dots + \lambda_n v_n = 0 ~ \implies ~ v_i = \left(-\dfrac{\lambda_1}{\lambda_i}\right)v_1 + \cdots + \left(-\dfrac{\lambda_{i - 1}}{\lambda_i}\right)v_{i - 1} + \left(-\dfrac{\lambda_{i + 1}}{\lambda_i}\right)v_{i + 1} + \cdots + \left(-\dfrac{\lambda_n}{\lambda_i}\right)v_n. \]
    In other words, because $\lambda_i$ is nonzero, $v_i$ must already be in the span of $S'$, meaning $\mathrm{span}(S') = \mathrm{span}(S)$.
    
    So as long as $S$ is not linearly independent, we can keep removing vectors from $S$ while preserving the property that $S$ spans $V$.  Since $S$ is finite, this process of removing vectors from $S$ must end eventually, at which point $S$ must be linearly independent and thus a basis of $V$. ~ $\blacksquare$
\end{proof}

\begin{remark}
It turns out that infinite-dimensional vector spaces have bases too.  But the proof is much more technical---in fact, for the purposes of this course, we'll disregard infinite-dimensional vector spaces for their excessive technicality and lack of mathematical beauty.
\end{remark}

Now, we'd really like to make the following definition.

\textbf{Definition? (Dimension?).} The dimension of a vector space $V$ is the size of any basis $S \subseteq V$\dots?

However, we need to check that this definition is self-consistent.  Must any two bases $S_1$ and $S_2$ of a vector space $V$ have the same size?  It turns out the answer is yes!  To prove this, we need the following theorem:

\begin{theorem}{Spans Greater Than Independents}
    If $S_1 = \{v_1, \dots, v_r\} \subseteq V$ spans $V$, and $S_2 = \{w_1, \dots, w_s\} \subseteq V$ is linearly independent, then $r \geq s$. 
\end{theorem}

\begin{proof}
    Argue by contradiction: suppose $r < s$.  Because $S_1$ spans $V$, we may express every single $w_j \in S_2$ as a linear combination of elements in $S_1$.  Suppose these linear combinations look like this:
    \[ \text{For all } j \in \{1, \dots, s\}, ~ \text{ we have } \sum_{i = 1}^r A_{ij} v_i = w_j. \]
    Note that, as portrayed above, the $\{A_{ij}\}$ are nothing more than plain coefficients of linear combinations.  But suppose we interpreted the $\{A_{ij}\}$ as entries of an $r \times s$ matrix $A$, anyway.
    
    Then if we perform row-reduction on $A$, the resulting matrix must have a column without a pivot, since $r < s$.  This means that there is some nonzero $x = \{x_1, \dots, x_s\}\in F^s$ such that $Ax = 0$.

    The finish is to violate the linear independence of $S_2$ by using $\{x_i\}$ as coefficients of a linear combination of $S_2$.
    \[ \sum_{j = 1}^s x_j w_j = \sum_{j = 1}^s x_j \left(\sum_{i = 1}^r A_{ij}v_i\right) = \sum_{i = 1}^r \left(\sum_{j = 1}^s A_{ij} x_j \right) v_i = \sum_{i = 1}^r (0)v_i = 0.  \]
    This contradicts the linear independence of $S_2$, as desired. ~ $\blacksquare$
\end{proof}

An immediate corollary of the above theorem is exactly what we want.

\begin{theorem}{Consistency of Basis Size}
    For any two bases $S_1 = \{v_1, \dots, v_r \} \subseteq V$ and $S_2 = \{w_1, \dots, w_s \} \subseteq V$ of a vector space $V$, it must be that $r = s$.
\end{theorem}

\begin{proof}
    Note that $S_1$ spans $V$ and $S_2$ is linearly independent, so $r \geq s$.  Also, $S_1$ is linearly independent and $S_2$ spans $V$, so $s \geq r$.  These two inequalities together imply $s = r$, done. ~ $\blacksquare$
\end{proof}

And so we can finally write the definition of dimension without question marks.

\textbf{Definition (Dimension).} Let $S$ be any basis of a vector space $V$.  Then the \textbf{dimension} of $V$ equals $|S|$.

\end{document}